\subsection*{迟到请求应回到哪笔成交}

重新从账户 1 的 10000 分、空购买记录开始。Alice 的 S1 使用 K1 买 \code{paper-001}，Q7 为 3000，事务提交 P1、O1、D1，当次余额为 7000。成功响应丢失，客户端仍把 S1 保存为待确认。此时原购买已经成立，手机界面尚未取得结果。

Alice 通过另一个明确确认动作发起 1000 分的合法购买，身份为 S2，使用 K2。为了只观察身份与重试交错，这里假定这项购买的商品、报价和权限都已经验证，按同一套事务规则完成 P2、O2、D2。它从 7000 扣到 6000，响应正常返回。S2 的成立无需等 S1 的客户端先获知结果，账户内可以同时存在一个已成交待确认操作和另一个已确认操作。

S1 的恢复请求 R2 随后到达，仍携带 E1、K1 和原输入。它对 S1 的唯一插入返回零行，下一条 SELECT 找到 P1，核对输入相同，返回保存的 P1 成功结果。它不会再次进入余额更新。数据库仍保留两项合法操作，当前余额为 \(10000-3000-1000=6000\)。

\begin{figure}[htbp]
\centering
\begin{tikzpicture}[x=1cm,y=1cm]
\node[booknode,text width=3.7cm] (p1) at (2,0) {S1 首次购买\\10000 扣 3000 → 7000};
\node[booknode,text width=3.7cm] (p2) at (7,0) {S2 独立购买\\7000 扣 1000 → 6000};
\node[bookresult,text width=3.7cm] (retry) at (12,0) {S1 迟到重试\\只读取 P1 原结果};
\draw[bookarrow] (p1)--(p2);\draw[bookarrow] (p2)--(retry);
\node[booknode,text width=3.7cm] (d1) at (2,-1.8) {P1、O1、D1\\当次余额 7000};
\node[booknode,text width=3.7cm] (d2) at (7,-1.8) {P2、O2、D2\\当前余额 6000};
\node[booknode,text width=3.7cm] (rr) at (12,-1.8) {返回 P1、O1、D1\\保留当次 7000};
\draw[bookarrow] (p1)--(d1);\draw[bookarrow] (p2)--(d2);\draw[bookarrow] (retry)--(rr);
\node[booknote,text width=13cm] at (7,-3.1) {原操作快照与当前账户分别保留时点；重试不覆盖 S2 的已提交变化};
\end{tikzpicture}
\caption{另一笔合法扣款插入以后，迟到重试回到原购买而不再更新余额。}
\label{fig:ch04-intervening-purchase}
\end{figure}


原结果里的 \code{balance\_after\_minor=7000} 描述 D1 当时扣款完成后的余额；账户查询里的 \code{current\_balance\_minor=6000} 描述查询时的余额。两个值分别有记录与时点支持。恢复请求如果把 P1 的当次余额改成 6000，同一条扣款流水的历史字段就会随其他交易改变。

\begin{BookCode}
{
  "purchase_id": "P1",
  "status": "succeeded",
  "order_id": "O1",
  "debit_id": "D1",
  "amount_minor": 3000,
  "balance_after_minor": 7000
}
\end{BookCode}

这段响应只列出 P1 原结果中与本节有关的字段；服务为 R2 另加请求编号和响应时间。页面若同时展示当前余额，就明确标记查询时点，并从账户入口读取 6000。恢复 P1 要找回这次购买得到了什么结论，查询账户则说明现在还能使用多少金额。

\subsection*{另一位观察者已经读到了什么}

Alice 的手机等待 P1 响应时，另一台已授权设备可能已经在历史中读到 O1 和 D1；它又完成 S2，读到账户为 6000。手机后来恢复 S1，应当取得能与这条历史共同成立的结果：P1 已经成功、当次余额为 7000，而 S2 仍然存在。把所有界面上的余额强行改成同一值，会混淆购买发生时的快照与当前账户状态。

另一台设备取得 P1 的购买编号以后，还可能继续打开 P1 详情。如果恢复代码删掉 P1，另建 P3 来替代，手机或许最终显示成功，另一台设备却会失去原查询入口。沿同一 S 返回 P1，既保留它已经看到的成交，也让后来取得结果的手机指向同一笔购买。请求的返回次序可以变化，已提交的对象关联继续保持。

再从流水检查这条历史。D1 保存 3000 和当次 7000，D2 保存 1000 和当次 6000；两条合计 4000，与当前余额 \(10000-4000=6000\) 相符。R2 只是读取 P1，没有新增扣款，也没有删除 D2。若它再次把余额写成 7000，就会覆盖 S2 的合法变化；下一节将把这条错误路径与数学幂等等式放在一起分析。

\subsection*{已有终态与尚无终态的分界}

成功结果与余额不足都需要这种恢复方式：已有终态时，后来的独立交易不改写原结论。余额不足以后充值的完整响应与新购买规则集中在第 5.4 节；当前例子先确定成功购买的历史字段与当前字段分别来自哪里。

原事务完整回滚则没有已提交终态。若 S1 回滚，S2 先从 10000 扣去 1000，随后 S1 恢复，它要重新竞争身份，并依据当时的 9000 执行条件扣款；成功后当次余额为 6000。这与 S1 先提交到 7000、S2 再提交到 6000 的最终余额相同，S1 的原结果却不同。区别来自 S1 在哪一次事务中真正提交，候选编号或先前的超时不能替代这项事实。

因此，恢复前需要同时沿完整身份读取终态与原结果。S1 已有终态，就回到它保存的编号和字段；S1 尚无终态，就继续受约束地尝试，并接受其他合法购买已经改变账户的事实。第五章将进一步给出完成记录中需要保存的具体字段，使这些区别在价格与接口版本变化以后仍可恢复。
\FloatBarrier
